\documentclass[12pt]{article}

\usepackage[T1]{fontenc}
\usepackage[utf8]{inputenc}
\usepackage{lmodern}
\usepackage{microtype}
\usepackage[letterpaper,margin=1in]{geometry}
\usepackage{amssymb,amsthm,mathtools}
\usepackage{titlesec}
\usepackage[hidelinks]{hyperref}

\hypersetup{
  pdftitle={Algebraic Independence in the Affine Case of Erdos Problem 270},
  pdfsubject={Algebraic independence for positive affine functions in Erdos Problem 270},
  pdfauthor={Costa Lambrinoudis}
}

\allowdisplaybreaks
\linespread{1.02}
\setlength{\parindent}{0pt}
\setlength{\parskip}{0.2\baselineskip}
\emergencystretch=2em
\titleformat{\section}{\raggedright\normalfont\Large\bfseries}{\thesection}{1em}{}
\titlespacing*{\section}{0pt}{3.5ex plus 1ex minus 0.2ex}{2.3ex plus 0.2ex}
\titleformat{\subsection}{\raggedright\normalfont\large\bfseries}{\thesubsection}{1em}{}
\titlespacing*{\subsection}{0pt}{3.25ex plus 1ex minus 0.2ex}{1.5ex plus 0.2ex}
\newtheorem{theorem}{Theorem}
\newtheorem{lemma}{Lemma}[section]
\newtheorem*{descentlemma}{Descent lemma}
\newtheorem*{corollary}{Corollary}

\begin{document}

\begin{center}
{\LARGE Algebraic Independence in the Affine Case of Erd\H{o}s Problem 270}\\[1.5ex]
{\large Costa Lambrinoudis}\\[0.5ex]
{23 August 2026}\\[1.75ex]
{\small\itshape AI Disclosure: This manuscript was written and checked using generative AI.}
\end{center}

\begin{quote}
\small
Erd\H{o}s Problem 270 asks whether the series of reciprocals of $(n+1)(n+2)\cdots(n+f(n))$ must be irrational whenever $f(n)\to\infty$. Crmari\'c and Kova\v{c} showed that the answer is negative without a regularity assumption, while the problem remains open for nondecreasing $f$. Erd\H{o}s and Graham singled out the simplest such choice, $f(n)=n$. We study the entire positive affine subfamily $f(n)=an+b$, whose members are
\[
C_{a,b}:=\sum_{n=1}^{\infty}\frac{n!}{((a+1)n+b)!}
\]
for $a\geq1$ and $b\geq1-a$. For each fixed $a$, we prove that $C_{a,0},\ldots,C_{a,a-1},C_{a,a+1}$ are algebraically independent over $\overline{\mathbb Q}$ and that every other admissible $C_{a,b}$ is a rational affine combination of them. Thus the fixed-slope family has transcendence degree $a+1$; in particular, every positive affine case of Problem 270 is transcendental. We begin with an elementary irrationality proof for $0\leq b\leq a$, which includes the Erd\H{o}s--Graham example. A Gaussian integral representation then gives a short transcendence proof for that example. For the general theorem, we pass to hypergeometric $E$-functions and combine algebraic-independence results of Salikhov and Salikhov--Viskina with a logarithmic obstruction at infinity and Beukers's specialization theorem.
\end{quote}

\section{Irrationality in the affine family}
\label{sec:irrationality}

Erd\H{o}s and Graham ended their chapter on irrationality and transcendence with the question now known as Erd\H{o}s Problem 270. Given a positive integer-valued function $f$ with $f(n)\to\infty$, they considered
\[
S(f):=\sum_{n=1}^{\infty}
\frac1{(n+1)(n+2)\cdots(n+f(n))},
\]
and asked whether $S(f)$ must be irrational. They expected an affirmative answer when $f$ is nondecreasing~\cite[p.~66]{erdos-graham}. Crmari\'c and Kova\v{c} later showed that some regularity assumption is necessary: as $f$ ranges over all such functions, $S(f)$ takes every value in $(0,\infty)$~\cite[Theorem~1]{crmaric-kovac}. The nondecreasing problem remains open~\cite{bloom-270}.

For integers $a\geq1$ and $b\geq0$, the affine choice $f(n)=an+b$ gives
\[
S(f)=\sum_{n=1}^{\infty}\frac1{(n+1)\cdots((a+1)n+b)}
=\sum_{n=1}^{\infty}\frac{n!}{((a+1)n+b)!}
=:C_{a,b}.
\]
Erd\H{o}s and Graham singled out the simplest choice, $f(n)=n$, and described the resulting series as difficult to treat~\cite[pp.~65--66]{erdos-graham}. It is the case $a=1$, $b=0$ of the family above. The following argument proves irrationality throughout the larger range $0\leq b\leq a$.

Fix $a\geq1$ and $0\leq b\leq a$, and put
\[
D_n:=\frac{((a+1)n+b)!}{n!}.
\]
Then
\[
\frac{D_{n+1}}{D_n}
=\frac{\displaystyle\prod_{j=1}^{a+1}\bigl((a+1)n+b+j\bigr)}{n+1}.
\]
The factor with $j=a+1-b$ is $(a+1)(n+1)$. Cancelling $n+1$ gives the positive integer
\[
r_n:=\frac{D_{n+1}}{D_n}
=(a+1)
\prod_{\substack{1\leq j\leq a+1\\j\neq a+1-b}}
\bigl((a+1)n+b+j\bigr).
\]
Thus $D_n\mid D_{n+1}$. Moreover, $r_n\geq a+1\geq2$, and every factor in the product defining $r_n$ increases with $n$. Hence $r_n$ is increasing and tends to infinity. Let
\[
P_N:=\sum_{n=1}^{N}\frac1{D_n},
\qquad
R_N:=C_{a,b}-P_N.
\]
Since $D_1\mid D_2\mid\cdots\mid D_N$, we have $D_NP_N\in\mathbb Z$. Moreover,
\[
\begin{aligned}
0<D_NR_N
&=\sum_{k=1}^{\infty}
\frac1{r_Nr_{N+1}\cdots r_{N+k-1}}\\
&\leq\sum_{k=1}^{\infty}\frac1{r_N^k}
=\frac1{r_N-1}
\longrightarrow0.
\end{aligned}
\]
If $C_{a,b}=p/q$, then $D_NR_N$ is positive and belongs to $q^{-1}\mathbb Z$, so it is at least $1/q$. This is impossible for large $N$. Hence $C_{a,b}$ is irrational for every $a\geq1$ and $0\leq b\leq a$.

\begin{corollary}
The Erd\H{o}s--Graham series
\[
C_{1,0}=\sum_{n=1}^{\infty}\frac{n!}{(2n)!}
\]
is irrational.
\end{corollary}

Crmari\'c and Kova\v{c} returned to this example and reported that they knew no proof of its irrationality~\cite[p.~55, following (1.1)]{crmaric-kovac}. After their paper appeared, they found the $a=1$, $b=0$ specialization of the argument above and suggested that similar reasoning might extend further; Kova\v{c} posted this observation on the Erd\H{o}s Problems forum in July 2026~\cite{kovac-forum}. We obtained the general argument independently before learning of their post.

\section{Transcendence of \texorpdfstring{$C_{1,0}$}{C(1,0)}}
\label{sec:a1-transcendence}

The same constant has a Gaussian integral representation. Crmari\'c and Kova\v{c} record the identity~\cite[p.~55 and footnote~1]{crmaric-kovac}
\[
C_{1,0}=e^{1/4}\int_0^{1/2}e^{-t^2}\,dt.
\tag{1}
\]
They obtain it by expanding $e^{1/4-t^2}$ and integrating termwise.

In a 2023 MathOverflow discussion of the neighboring integral $\int_0^1e^{-t^2}\,dt$, Weber and Pinelis~\cite{mathoverflow} considered the $E$-function
\[
y(x):=\frac1{\sqrt{x}}\int_0^{\sqrt{x}}e^{-t^2}\,dt
=\sum_{n=0}^{\infty}\frac{(-1)^n}{2n+1}\frac{x^n}{n!}
\]
and used the Siegel--Shidlovsky theorem to prove that $y(1)$ and $e^{-1}$ are algebraically independent. The differential system in their argument is ordinary at every nonzero algebraic point, so applying the same theorem at $x=1/4$ gives the algebraic independence of $y(1/4)$ and $e^{-1/4}$. Since
\[
y(1/4)=2\int_0^{1/2}e^{-t^2}\,dt,
\qquad
C_{1,0}=\frac{y(1/4)}{2e^{-1/4}},
\]
an algebraic value of $C_{1,0}$ would give a nontrivial algebraic relation between $y(1/4)$ and $e^{-1/4}$. Hence $C_{1,0}$ is transcendental. The MathOverflow discussion did not mention either the factorial series or this specialization.

We now turn to the hypergeometric $E$-functions associated with the full affine family.

\section{Algebraic independence in the affine family}
\label{sec:affine-family}

\begin{theorem}
\label{thm:main}
For every integer $a\geq1$, the numbers
\[
C_{a,0},C_{a,1},\ldots,C_{a,a-1},C_{a,a+1}
\]
are algebraically independent over $\overline{\mathbb Q}$. Every $C_{a,b}$ with $b\geq1-a$ is a rational affine combination of these $a+1$ numbers, and
\[
\operatorname{trdeg}_{\overline{\mathbb Q}}
\overline{\mathbb Q}(C_{a,b}:b\geq1-a)=a+1.
\]
Moreover, every $C_{a,b}$ in this range is transcendental.
\end{theorem}

We first prove the transcendence assertion for every admissible intercept. We then use the same functions to determine all algebraic relations at each fixed slope.

For all integers $a\geq1$ and $b\geq0$, define
\[
F_{a,b}(z):=b!\sum_{n=0}^{\infty}
\frac{n!}{((a+1)n+b)!}z^{an}.
\tag{2}
\]
Then
\[
F_{a,b}(1)=1+b!C_{a,b}.
\tag{3}
\]
For $b\geq0$, it is therefore enough to prove that $F_{a,b}(1)$ is transcendental. Beukers's theorem reduces this to finding a differential system, ordinary at $1$, whose components include $F_{a,b}$ and are linearly independent over $\overline{\mathbb Q}(z)$.

\subsection{The range \texorpdfstring{$0\leq r\leq a$}{0 <= r <= a}}
\label{sec:initial-range}

Assume first that $0\leq r\leq a$. Put
\[
J_{a,r}:=\{1,\ldots,a+1\}\setminus\{a+1-r\},
\qquad
\beta_j:=\frac{r+j}{a+1}\quad(j\in J_{a,r}).
\]
The omitted parameter is the unique one equal to $1$. The multiplication formula gives
\[
F_{a,r}(z)
={}_1F_a\!\left(
\begin{matrix}1\\(\beta_j)_{j\in J_{a,r}}\end{matrix};
\frac{z^a}{(a+1)^{a+1}}
\right).
\tag{4}
\]
Indeed,
\[
\prod_{j=1}^{a+1}
\left(\frac{r+j}{a+1}\right)_n
=\frac{((a+1)n+r)!}{r!(a+1)^{(a+1)n}},
\]
and deleting the factor $(1)_n=n!$ yields~(4).

To match the notation of the cited results, set
\[
\rho_a:=\frac{a}{(a+1)^{(a+1)/a}},
\qquad
\lambda_j:=\beta_j-1,
\]
and define
\[
\varphi(w):=\sum_{n=0}^{\infty}
\frac{(w/a)^{an}}
{\displaystyle\prod_{j\in J_{a,r}}(\lambda_j+1)_n}.
\]
Then $F_{a,r}(z)=\varphi(\rho_a z)$. This is the $l=0$, $t=a$ case of Salikhov's formula~\cite[eq.~(1)]{salikhov}. Here $l=0$ means that Salikhov's numerator contains no Pochhammer factors; in standard hypergeometric notation the parameter $1$ in~(4) remains because $(1)_n/n!=1$.

Let $\theta=z\,d/dz$ and define
\[
W_{a,r}:=\operatorname{span}_{\mathbb C(z)}
\{1,F_{a,r},\theta F_{a,r},\ldots,\theta^{a-1}F_{a,r}\}.
\tag{5}
\]
The key functional statement is
\[
\dim_{\mathbb C(z)}W_{a,r}=a+1.
\tag{6}
\]
For odd $a$, Salikhov's Theorem~1 applies to this $l=0$ function~\cite[Theorem~1]{salikhov}. The parameters are rational, and none is a negative integer. The remaining hypothesis requires that their multiset modulo integers not be invariant under translation by $1/d$ for any divisor $d>1$ of $a$. In the present notation, this multiset is
\[
S_a=\left\{\frac1{a+1},\ldots,\frac a{a+1}\right\}\subset\mathbb Q/\mathbb Z.
\]
If $d>1$ divides $a$, then $\gcd(d,a+1)=1$, so $1/d$ does not lie in the subgroup $(a+1)^{-1}\mathbb Z/\mathbb Z$ containing $S_a$. Translation by $1/d$ therefore cannot preserve $S_a$, which verifies Salikhov's parameter condition. The theorem gives algebraic independence of $\varphi(\alpha),\varphi'(\alpha),\ldots,\varphi^{(a-1)}(\alpha)$ for every nonzero algebraic $\alpha$.

When $a$ is even, the theorem of Salikhov--Viskina applies with $t=a$. If the $\lambda_j$ are rational and not negative integers, it states that for every nonzero algebraic $\alpha$ the numbers $1,\varphi(\alpha),\delta\varphi(\alpha),\ldots,\delta^{a-1}\varphi(\alpha)$ are linearly independent over $\overline{\mathbb Q}$, where $\delta=w\,d/dw$ is their Euler operator~\cite[p.~833]{salikhov-viskina}. Unlike Salikhov's odd-order result, this gives linear rather than algebraic independence; it includes the constant $1$ and imposes no translation condition.

In either parity, the cited theorem implies that the values of $1,F_{a,r},\theta F_{a,r},\ldots,\theta^{a-1}F_{a,r}$ are linearly independent over $\overline{\mathbb Q}$ at every nonzero algebraic point. At such a point $\alpha$, the Euler derivatives $(\theta^kF_{a,r})(\alpha)$ and the ordinary derivatives $F_{a,r}^{(k)}(\alpha)$ are related by a triangular algebraic matrix with diagonal entries $1,\alpha,\ldots,\alpha^{a-1}$. The two forms of the independence statement are therefore equivalent.

\begin{descentlemma}
Let $g_1,\ldots,g_m\in\overline{\mathbb Q}[[z]]$. If these series are linearly dependent over $\mathbb C(z)$, then they are linearly dependent over $\overline{\mathbb Q}(z)$.
\end{descentlemma}

\begin{proof}
Clear the denominators in a nonzero relation and let $D$ be the largest degree of its polynomial coefficients. Writing those coefficients as $\sum_{d=0}^{D}c_{i,d}z^d$ and comparing Taylor coefficients gives a homogeneous linear system over $\overline{\mathbb Q}$ in the finitely many unknowns $c_{i,d}$. Its row space has a finite basis, so the system has the same kernel as a finite matrix over $\overline{\mathbb Q}$. If its kernel over $\mathbb C$ is nonzero, then its kernel over $\overline{\mathbb Q}$ is nonzero, and the corresponding polynomial coefficients give the required relation.
\end{proof}

Every function in~(5) has algebraic Taylor coefficients, so the descent lemma implies~(6). Otherwise a relation over $\mathbb C(z)$ would descend to $\overline{\mathbb Q}(z)$. Specializing at a nonzero algebraic point where its coefficients are defined and not all zero would contradict the cited value independence.

Applying the value-independence statement at $z=1$ proves that $F_{a,r}(1)$, and hence $C_{a,r}$ by~(3), is transcendental for $0\leq r\leq a$.

The same statement handles affine functions with negative intercept that remain positive for every $n\geq1$. Suppose that $1-a\leq b<0$ and put $r=a+1+b$, so that $a\geq2$ and $2\leq r\leq a$. Shifting the summation index by writing $n=m+1$ gives
\[
\begin{aligned}
C_{a,b}
&=\sum_{m=0}^{\infty}\frac{(m+1)!}{((a+1)m+r)!}\\
&=\frac1{r!}\left(F_{a,r}(1)+\frac1a\theta F_{a,r}(1)\right).
\end{aligned}
\]
If $C_{a,b}$ were algebraic, this identity would give a nontrivial linear relation over $\overline{\mathbb Q}$ among $1$, $F_{a,r}(1)$, and $\theta F_{a,r}(1)$. Hence $C_{a,b}$ is transcendental for $1-a\leq b<0$ as well.

To treat $b>a$, we compare the formal solutions of the relevant differential equations at infinity. The generalized hypergeometric equation for~(4) can be written as
\[
\left[
\prod_{j\in J_{a,r}}
\left(\frac{\theta}{a}+\beta_j-1\right)
-\frac{z^a}{(a+1)^{a+1}}
\right]F_{a,r}
=\prod_{j\in J_{a,r}}(\beta_j-1).
\tag{7}
\]
The right side is nonzero because the parameter equal to $1$ was removed. To describe the formal system at infinity, put
\[
x:=z^{-1},\qquad K:=\mathbb C((x)),\qquad \kappa_a:=(a+1)^{-(a+1)},
\]
so that $\theta=-x\,d/dx$ on $K$, and write
\[
c_{a,r}:=\prod_{j\in J_{a,r}}(\beta_j-1),\qquad P_{a,r}(u):=\prod_{j\in J_{a,r}}(u+\beta_j-1)=\sum_{j=0}^{a}p_ju^j,\qquad p_a=1.
\]
Let
\[
Y_0:=1,\qquad Y_j:=\theta^{j-1}F_{a,r}\quad(1\leq j\leq a).
\]
Equation~(7) is equivalent to the augmented companion system
\[
\theta Y_0=0,\qquad \theta Y_j=Y_{j+1}\quad(1\leq j<a),
\]
\[
\theta Y_a=a^ac_{a,r}Y_0+a^a(\kappa_a z^a-p_0)Y_1-\sum_{j=1}^{a-1}a^{a-j}p_jY_{j+1}.
\]
After expanding rational functions at infinity, this system defines a rank-$(a+1)$ differential module $\mathcal M_{a,r}$ over $K$.

\begin{lemma}[Zero-exponential formal component]
\label{lem:zero-component}
The Levelt--Turrittin decomposition of $\mathcal M_{a,r}$ is
\[
\mathcal M_{a,r}=\mathcal M_0\oplus\bigoplus_{\zeta^a=1}\mathcal M_\zeta,
\]
where $\mathcal M_0$ has rank one and exponential part $0$, while each $\mathcal M_\zeta$ has rank one and exponential part $\rho_a\zeta z$. The zero summand has the formal solution vector
\[
\widehat Y^{(0)}:=\left(1,\widehat F_{a,r}^{(0)},\theta\widehat F_{a,r}^{(0)},\ldots,\theta^{a-1}\widehat F_{a,r}^{(0)}\right)^{\mathsf T},
\]
where
\[
\widehat F_{a,r}^{(0)}(z)\in z^{-a}\mathbb C[[z^{-a}]].
\tag{8}
\]
More explicitly, if
\[
G=r_0+\sum_{k=0}^{a-1}r_{k+1}\theta^kF_{a,r}\in W_{a,r},\qquad r_0,\ldots,r_a\in\mathbb C(z),
\]
then its zero-exponential projection is
\[
\pi_0(G)=r_0+\sum_{k=0}^{a-1}r_{k+1}\theta^k\widehat F_{a,r}^{(0)}\in K.
\]
The projection extends $K$-linearly to $K\otimes_{\mathbb C(z)}W_{a,r}$ and commutes with $\theta$. In particular, $\pi_0(G)$ is a Laurent series without powers of $\log z$.
\end{lemma}

\begin{proof}
Writing $\widehat F_{a,r}^{(0)}=\sum_{m\geq1}h_mz^{-am}$ in~(7) gives
\[
h_1=-\frac{c_{a,r}}{\kappa_a},
\qquad
h_{m+1}=\frac{P_{a,r}(-m)}{\kappa_a}h_m
\quad(m\geq1).
\tag{9}
\]
This recurrence constructs the solution~(8). It is unique among Laurent-series particular solutions. Indeed, the difference of two such solutions would be a nonzero Laurent-series solution of the homogeneous equation. If $dz^\nu$ were its leading monomial, the term $-\kappa_a z^a dz^\nu$ could not be cancelled by $P_{a,r}(\theta/a)$, which does not raise exponents. Consequently $\widehat Y^{(0)}$ is a nonzero formal solution of the augmented system with exponential part $0$.

The homogeneous equation associated with~(7) has characteristic equation
\[
\lambda^a=\frac{a^a}{(a+1)^{a+1}},
\]
so its $a$ characteristic numbers are the distinct nonzero numbers $\rho_a\zeta$. Its formal homogeneous solutions have the form
\[
e^{\rho_a\zeta z}z^{\mu_\zeta}\mathbb C[[z^{-1}]],
\qquad \zeta^a=1,
\]
and are log-free because every characteristic number is simple; see Salikhov~\cite[Remark~2, p.~206]{salikhov}. Each of the $a$ distinct nonzero exponential parts therefore contributes a rank-one summand.

The Levelt--Turrittin theorem decomposes $\mathcal M_{a,r}$ as a direct sum of differential modules indexed by their exponential parts~\cite[Chapter~3, Theorem~3.1]{vdput-singer}. No ramification is needed because $0$ and the polynomials $\rho_a\zeta z$ lie over $K$. The nonzero exponential summands account for rank $a$, and $\widehat Y^{(0)}$ supplies a nonzero solution in the zero-exponential summand. Since $\mathcal M_{a,r}$ has rank $a+1$, the zero summand has rank exactly one, and its formal solution space is spanned by $\widehat Y^{(0)}$.

By~(6), the displayed expression for $G$ in the statement is unique. Expanding the coefficients $r_j$ in $K$ and projecting the companion vector onto the line spanned by $\widehat Y^{(0)}$ gives the stated formula for $\pi_0(G)$. Because the Levelt--Turrittin summands are differential modules, this projection is $K$-linear and commutes with $\theta$. Every entry of $\widehat Y^{(0)}$ lies in $K$, so the formula also shows directly that $\pi_0(G)$ contains no logarithms.
\end{proof}

\subsection{Euler reduction for \texorpdfstring{$b>a$}{b > a}}
\label{sec:euler-reduction}

Suppose that $b>a$ and write
\[
b=q(a+1)+r,
\qquad q\geq1,
\qquad 0\leq r\leq a.
\tag{10}
\]
Put
\[
\beta_j:=\frac{b+j}{a+1},
\qquad 1\leq j\leq a+1,
\qquad
x:=\frac{z^a}{(a+1)^{a+1}}.
\]
None of the lower parameters is equal to $1$, and the multiplication formula gives
\[
F_{a,b}(z)
={}_2F_{a+1}\!\left(
\begin{matrix}1,1\\\beta_1,\ldots,\beta_{a+1}\end{matrix};
x
\right).
\tag{11}
\]
The function $F_{a,b}$ in~(11) is a hypergeometric $E$-function: its parameters are positive rational numbers, and with two upper and $a+1$ lower parameters, the required pullback degree is $(a+1)-2+1=a$. Writing its differential equation in terms of $\theta_x=x\,d/dx$, factoring off $\theta_x$, and evaluating at $x=0$ to determine the resulting constant gives
\[
\left[
\prod_{j=1}^{a+1}\left(\frac{\theta}{a}+\beta_j-1\right)
-\frac{z^a}{(a+1)^{a+1}}\left(\frac{\theta}{a}+1\right)
\right]F_{a,b}
=\prod_{j=1}^{a+1}(\beta_j-1).
\tag{12}
\]
This is an order-$(a+1)$ inhomogeneous equation. Its leading ordinary-derivative coefficient is a nonzero constant multiple of $z^{a+1}$, so its companion system has no nonzero finite singularity.

Coefficient comparison in~(2) gives the Euler reduction
\[
\boxed{
\prod_{k=1}^{q}\left(\frac{\theta}{a}+k\right)F_{a,b}(z)
=\frac{b!}{r!}z^{-aq}
\left[
F_{a,r}(z)
-r!\sum_{m=0}^{q-1}
\frac{m!}{((a+1)m+r)!}z^{am}
\right].}
\tag{13}
\]
The operator on the left multiplies the coefficient of $z^{an}$ by
\[
\prod_{k=1}^{q}(n+k)=\frac{(n+q)!}{n!}.
\]
The index shift $n\mapsto n+q$ gives the right side of~(13). In particular, the differential module generated by $1$ and $F_{a,b}$ contains $W_{a,r}$.

\subsection{Resonance}
\label{sec:resonance}

Assume for contradiction that $F_{a,b}\in W_{a,r}$ and apply the projection $\pi_0$ of Lemma~\ref{lem:zero-component} to~(13). This projection commutes with the Euler operator on the left. On the right, the contribution of $z^{-aq}F_{a,r}$ begins at $z^{-a(q+1)}$ by~(8), whereas the highest term of the polynomial subtracted in~(13), corresponding to $m=q-1$, contributes the nonzero $z^{-a}$ coefficient
\[
-\frac{b!(q-1)!}{((a+1)(q-1)+r)!}.
\tag{14}
\]

Under this assumption, the lemma shows that the zero-exponential component of $F_{a,b}$ is a Laurent series $L(z)\in\mathbb C((z^{-1}))$ without logarithms. The Euler operator on the left of~(13) acts diagonally on monomials, and its multiplier on $z^{-a}$ is
\[
\prod_{k=1}^{q}(-1+k)=0.
\tag{15}
\]
It follows that the $z^{-a}$ coefficient in the zero-exponential component of the left side of~(13) must vanish, contradicting~(14). Therefore
\[
F_{a,b}\notin W_{a,r}.
\tag{16}
\]

The projected equation therefore requires a logarithmic term. Since
\[
\prod_{k=1}^{q}\left(\frac{\theta}{a}+k\right)
\bigl(z^{-a}\log z\bigr)
=\frac{(q-1)!}{a}z^{-a},
\]
comparison with~(14) gives
\[
\widehat F_{a,b}^{(0)}(z)
=-\frac{a\,b!}{((a+1)(q-1)+r)!}\,z^{-a}\log z
+\cdots,
\tag{17}
\]
where the omitted terms do not affect the coefficient of $z^{-a}\log z$.

Define
\[
V_{a,b}:=\operatorname{span}_{\mathbb C(z)}
\{1,F_{a,b},\theta F_{a,b},\ldots,\theta^aF_{a,b}\}.
\tag{18}
\]
Equation~(12) makes $V_{a,b}$ stable under $\theta$ and gives $\dim V_{a,b}\leq a+2$. Equation~(13) gives $W_{a,r}\subseteq V_{a,b}$. By~(6) and~(16), the additional element $F_{a,b}$ is not in $W_{a,r}$, so
\[
\dim V_{a,b}\geq(a+1)+1=a+2.
\]
Hence $\dim V_{a,b}=a+2$, and
\[
1,F_{a,b},F_{a,b}',\ldots,F_{a,b}^{(a)}
\tag{19}
\]
are linearly independent over $\mathbb C(z)$. Thus~(12) has minimal possible order among inhomogeneous equations over $\mathbb C(z)$ satisfied by $F_{a,b}$.

\subsection{Values at \texorpdfstring{$z=1$}{z=1}}
\label{sec:evaluation}

By~(11), $F_{a,b}$ is an $E$-function, as are its derivatives. The functions in~(19) are linearly independent over $\overline{\mathbb Q}(z)$, and~(12) supplies a companion system with no nonzero finite singularity. Beukers~\cite[Corollary~1.4]{beukers} therefore gives the linear independence of their values at the ordinary algebraic point $z=1$. In particular, $F_{a,b}(1)$ is transcendental, and~(3) gives the transcendence of $C_{a,b}$.

Combining this with the cases $1-a\leq b\leq a$ treated above proves that $S(f)$ is transcendental for every positive integer-valued affine function $f(n)=an+b$ with $a\geq1$. It remains to determine the algebraic relations among the values with the same slope.

\subsection{Algebraic independence at fixed slope}
\label{sec:algebraic-independence}

Christopher D.~Long suggested extending the transcendence argument to algebraic independence~\cite{long-x}. For a fixed slope, the entire family cannot be algebraically independent because its members satisfy affine recurrences. We will prove that the maximal independent subfamily is
\[
C_{a,0},C_{a,1},\ldots,C_{a,a-1},C_{a,a+1}
\tag{20}
\]
and that
\[
\operatorname{trdeg}_{\overline{\mathbb Q}}
\overline{\mathbb Q}(C_{a,b}:b\geq1-a)=a+1.
\tag{21}
\]

Fix $a\geq1$, continue to write $\theta=z\,d/dz$, and abbreviate $F_b=F_{a,b}$. Direct comparison of the coefficients in~(2) gives
\[
F_{b-1}=F_b+\frac{a+1}{ab}\theta F_b
\qquad(1\leq b\leq a).
\tag{22}
\]
Indeed, the coefficient of $z^{an}$ on the right is
\[
\frac{b!n!}{((a+1)n+b)!}
\left(1+\frac{(a+1)n}{b}\right)
=\frac{(b-1)!n!}{((a+1)n+b-1)!}.
\]
Put $H=F_{a-1}$. Iterating~(22) gives
\[
F_j=P_j(\theta)H,
\qquad
P_j(X)=\prod_{k=j+1}^{a-1}
\left(1+\frac{a+1}{ak}X\right),
\qquad 0\leq j\leq a-1,
\tag{23}
\]
where an empty product is $1$. The polynomials $P_{a-1},P_{a-2},\ldots,P_0$ have degrees $0,1,\ldots,a-1$ and nonzero leading coefficients. Thus~(23) is an invertible triangular linear transformation between
\[
F_0,F_1,\ldots,F_{a-1}
\quad\text{and}\quad
H,\theta H,\ldots,\theta^{a-1}H.
\tag{24}
\]

For $H=F_{a-1}$, the lower shifts in Salikhov's notation are
\[
\lambda_j=\frac{j-2}{a+1},
\qquad 1\leq j\leq a+1,\quad j\neq2.
\tag{25}
\]
When $a$ is odd, Salikhov's Theorem~1 and the parameter check in Section~\ref{sec:initial-range} show that $H,\theta H,\ldots,\theta^{a-1}H$ are algebraically independent over $\mathbb C(z)$. When $a$ is even, Salikhov--Viskina show that algebraic dependence of the relevant function and its first $a-1$ Euler derivatives requires one of the shifts $\lambda_j$ to be an integer~\cite[p.~833]{salikhov-viskina}. None of the numbers in~(25) is an integer, so the same algebraic independence holds for even $a$. It follows from~(24) that
\[
F_0,F_1,\ldots,F_{a-1}
\tag{26}
\]
are algebraically independent over $\mathbb C(z)$ for every $a\geq1$.

It remains to adjoin one function from beyond the initial range. Let
\[
G:=F_{a+1},
\qquad
\mathcal K_a:=\mathbb C(z)(H,\theta H,\ldots,\theta^{a-1}H).
\]
The case $q=1$, $r=0$ of~(13) is
\[
\left(\frac{\theta}{a}+1\right)G
=(a+1)!z^{-a}(F_0-1).
\tag{27}
\]
The field $\mathcal K_a$ is stable under differentiation by the order-$a$ equation for $H$, and~(26) gives $\operatorname{trdeg}_{\mathbb C(z)}\mathcal K_a=a$.

\begin{lemma}[Logarithmic boundary obstruction]
\label{lem:field-log-obstruction}
Every element of $\mathcal K_a$ has a log-free formal expansion at infinity, whereas the zero-exponential component of every formal solution $G$ of~(27) contains
\[
-a(a+1)!z^{-a}\log z.
\tag{28}
\]
In particular, $G\notin\mathcal K_a$.
\end{lemma}

\begin{proof}
Apply Lemma~\ref{lem:zero-component} with $r=a-1$. A formal fundamental matrix for the augmented equation of $H$ consists of its rank-one zero-exponential solution and the $a$ rank-one solutions with distinct exponential parts $\rho_a\zeta z$. All of these solutions are log-free. Let $\mathcal P_a$ be the differential field generated over $\mathbb C((z^{-1}))$ by the entries of this matrix inside the standard formal Picard--Vessiot extension. It is generated by Laurent series, formal powers of $z$, and the exponentials $e^{\rho_a\zeta z}$, without adjoining a formal logarithm. The explicit admissible-field construction in Salikhov~\cite[pp.~207--210]{salikhov} treats the logarithm as a separate transcendental generator. Hence addition, multiplication, inversion, and differentiation in $\mathcal P_a$ do not introduce a power of $\log z$.

The formal expansions of $H,\theta H,\ldots,\theta^{a-1}H$ lie in $\mathcal P_a$, so the formal image of $\mathcal K_a$ is contained in $\mathcal P_a$ and is log-free. Project the right side of~(27) onto its zero-exponential component. The zero component of $F_0-1$ has constant term $-1$, and therefore the right side has leading term $-(a+1)!z^{-a}$. The operator $\theta/a+1$ annihilates $z^{-a}$. Since
\[
\left(\frac{\theta}{a}+1\right)
\bigl(z^{-a}\log z\bigr)=\frac1a z^{-a},
\]
the projected equation forces the term~(28). A solution of the homogeneous equation is a constant multiple of $z^{-a}$ and cannot remove the logarithm. Thus $G$ cannot belong to the log-free field $\mathcal K_a$.
\end{proof}

The lemma excludes rational membership in $\mathcal K_a$. The next elementary differential-field argument excludes algebraic membership.

\begin{lemma}[Trace descent]
\label{lem:trace-descent}
Let $K$ be a differential field of characteristic zero with algebraically closed field of constants, and let $R,p\in K$. Suppose that $y$ is algebraic over $K$ and satisfies $Dy+py=R$. Then the equation has a solution $y_0\in K$. If every solution of $Du+pu=0$ in an algebraic extension of $K$ already belongs to $K$, then $y\in K$.
\end{lemma}

\begin{proof}
Put $L=K(y)$ and $d=[L:K]$. Derivation commutes with the field trace, so
\[
y_0:=\frac1d\operatorname{Tr}_{L/K}(y)
\]
belongs to $K$ and satisfies $Dy_0+py_0=R$. The difference $y-y_0$ satisfies the homogeneous equation. The constants of the finite algebraic differential extension $L/K$ are algebraic over the constants of $K$ and hence introduce no new constants because the latter field is algebraically closed. The final assertion follows.
\end{proof}

Apply Lemma~\ref{lem:trace-descent} to~(27), with $D=\theta/a$ and $p=1$. The homogeneous solutions are $cz^{-a}$, with $c\in\mathbb C$, and already belong to $\mathcal K_a$. If $G$ were algebraic over $\mathcal K_a$, the trace lemma would therefore imply $G\in\mathcal K_a$, contrary to Lemma~\ref{lem:field-log-obstruction}. Hence $G$ is transcendental over $\mathcal K_a$, and the $a+1$ functions
\[
F_0,F_1,\ldots,F_{a-1},G
\tag{29}
\]
are algebraically independent over $\mathbb C(z)$.

These functions and the derivatives needed to close their equations form a system of $E$-functions over $\overline{\mathbb Q}(z)$ with no nonzero finite singularity. Its function field has transcendence degree $a+1$: the derivatives of $H$ lie in $\mathcal K_a$, while~(27) puts every derivative of $G$ in $\mathcal K_a(G)$. Beukers's Theorem~1.1~\cite[Theorem~1.1]{beukers} preserves this transcendence degree at the ordinary point $z=1$. By the triangular transformation~(24) and equation~(27), the corresponding value field is generated by
\[
F_0(1),F_1(1),\ldots,F_{a-1}(1),G(1).
\tag{30}
\]
The values in~(30) are therefore algebraically independent over $\overline{\mathbb Q}$. Equation~(3) gives the algebraic independence of the constants in~(20).

It remains to identify the rest of the fixed-slope family. Shifting the summation index gives
\[
C_{a,a}=(a+1)C_{a,0}-\frac1{a!}
\tag{31}
\]
and, for $b>a+1$,
\[
C_{a,b}=
\frac{C_{a,b-1}-(a+1)C_{a,b-a-1}+(a+1)/b!}{b-a-1}.
\tag{32}
\]
Thus every $C_{a,b}$ with $b\geq0$ is a rational affine combination of~(20). For $1-a\leq b<0$, set $r=a+1+b$. The index shift used in Section~\ref{sec:initial-range}, followed by~(22), gives
\[
C_{a,b}
=\frac1{r!}
+\frac{a+1-r}{a+1}C_{a,r}
+\frac1{a+1}C_{a,r-1}.
\tag{33}
\]
Here $2\leq r\leq a$, and~(31) eliminates $C_{a,a}$ when $r=a$. This proves the affine-span assertion in Theorem~\ref{thm:main}. The algebraic independence of~(20) supplies the matching lower bound in~(21) and completes the proof. In particular, the affine recurrences account for all algebraic relations at fixed slope: after any finite collection of the constants is expressed in the basis~(20), its ideal of polynomial relations is generated by the resulting affine linear relations.

\clearpage
\begin{thebibliography}{9}
\footnotesize
\raggedright
\setlength{\itemsep}{0pt}
\setlength{\parskip}{0pt}

\bibitem{erdos-graham}
P.~Erd\H{o}s and R.~L.~Graham, \emph{Old and New Problems and Results in Combinatorial Number Theory}, Monographies de L'Enseignement Math\'ematique \textbf{28} (1980), pp.~65--66.

\bibitem{bloom-270}
T.~F.~Bloom, Erd\H{o}s Problem \#270, \emph{Erd\H{o}s Problems}. \href{https://www.erdosproblems.com/270}{erdosproblems.com/270}; accessed 22 August 2026.

\bibitem{crmaric-kovac}
T.~Crmari\'c and V.~Kova\v{c}, On the irrationality of certain super-polynomially decaying series, \emph{Colloquium Mathematicum} \textbf{179} (2025), 55--68. \href{https://doi.org/10.4064/cm9628-5-2025}{doi:10.4064/cm9628-5-2025}.

\bibitem{kovac-forum}
V.~Kova\v{c}, comment on Erd\H{o}s Problem \#270, \emph{Erd\H{o}s Problems} forum, 12 July 2026. \href{https://www.erdosproblems.com/forum/thread/270\#post-7469}{Post 7469}; accessed 22 August 2026.

\bibitem{long-x}
C.~D.~Long, ``You can probably get algebraic independence here with a minor extension,'' post on X, 23 August 2026. \href{https://x.com/octonion/status/2091555025197662264}{Online post}; accessed 23 August 2026.

\bibitem{mathoverflow}
D.~Weber and I.~Pinelis, answers to ``Is the integral of $e^{-x^2}$ from $0$ to $1$ known to be irrational?'', MathOverflow, 26 December 2023. \href{https://mathoverflow.net/questions/461028/}{Online discussion}.

\bibitem{beukers}
F.~Beukers, A refined version of the Siegel--Shidlovskii theorem, \emph{Annals of Mathematics} \textbf{163} (2006), 369--379. \href{https://doi.org/10.4007/annals.2006.163.369}{doi:10.4007/annals.2006.163.369}.

\bibitem{salikhov}
V.~Kh.~Salikhov, A criterion for the algebraic independence of values of a class of hypergeometric $E$-functions, \emph{Mathematics of the USSR-Sbornik} \textbf{69} (1991), 203--226. \href{https://doi.org/10.1070/SM1991v069n01ABEH002113}{doi:10.1070/SM1991v069n01ABEH002113}.

\bibitem{salikhov-viskina}
V.~Kh.~Salikhov and G.~G.~Viskina, Algebraic relations between the hypergeometric $E$-function and its derivatives, \emph{Matematicheskie Zametki} \textbf{71} (2002), 832--844; English translation in \emph{Mathematical Notes} \textbf{71} (2002), 761--772. \href{https://doi.org/10.1023/A:1015864711107}{doi:10.1023/A:1015864711107}.

\bibitem{vdput-singer}
M.~van der Put and M.~F.~Singer, \emph{Galois Theory of Linear Differential Equations}, Grundlehren der mathematischen Wissenschaften \textbf{328}, Springer, 2003, Chapter~3, Theorem~3.1. \href{https://doi.org/10.1007/978-3-642-55750-7}{doi:10.1007/978-3-642-55750-7}.

\end{thebibliography}

\end{document}
